% =====================================================================
%  MỞ ĐẦU
% =====================================================================
\chapter*{Mở đầu}
\addcontentsline{toc}{chapter}{Mở đầu}
\markboth{Mở đầu}{Mở đầu}

Một cửa hàng thương mại điện tử phải quyết định giá mỗi ngày trong khi nhu cầu không chắc chắn và đối thủ cũng đang
quyết định giá. Có hai cách xây dựng ``bộ não'' cho quyết định đó. Cách thứ nhất là viết luật: người thiết kế đoán trước
các tình huống và ghi sẵn phản ứng. Cách thứ hai là cho tác tử tự học từ phần thưởng bằng học tăng cường (Reinforcement
Learning -- RL). Luật dễ hiểu và an toàn nhưng cứng; RL thích nghi được nhưng tốn dữ liệu, phụ thuộc seed và có thể học
sai nếu phần thưởng hoặc trạng thái thiết kế không khớp mục tiêu.

Bài tập Agent 02 đặt hai cách này cạnh nhau trong một hệ đa tác tử (Multi-Agent System -- MAS) tối giản: hai cửa hàng,
mỗi bước chọn giảm giá hoặc giữ giá~\cite{lab02}. Phần lý thuyết dựa trên ba tài liệu được giao: tutorial phát triển hệ đa
tác tử bằng Python và PyTorch~\cite{ch1:tutorial}, bài tổng quan học tăng cường của Naeem và cộng sự~\cite{ch2:naeem}, và
khảo sát cơ chế cộng tác của hệ đa tác tử dựa trên mô hình ngôn ngữ lớn của Tran và cộng sự~\cite{ch3:tran}.

\section*{Mục tiêu}
\begin{enumerate}
  \item Trình bày cơ sở thiết kế hệ đa tác tử, học tăng cường và cơ chế cộng tác đa tác tử theo ba tài liệu.
  \item Mô hình hóa bài toán theo tác tử, môi trường, trạng thái, hành động, phần thưởng, chính sách; tính sẵn chính sách
    tối ưu làm mốc.
  \item Cài đặt môi trường, tác tử Q-learning và tác tử luật bằng Python; tách huấn luyện và đánh giá.
  \item Thực hiện năm thí nghiệm A--E của đề bài và ba thí nghiệm mở rộng F--H; trả lời mười câu hỏi bằng số liệu đo được,
    cài đặt một thay đổi môi trường (câu 9) và một tác tử tốt hơn (câu 10).
  \item Hiện thực hệ đa tác tử hỗ trợ mua sắm theo tutorial, với bộ phân loại ý định PyTorch huấn luyện trên GPU.
\end{enumerate}

\section*{Cách đọc báo cáo}
Chương~\ref{chap:1}--\ref{chap:3} là lý thuyết. Chương~\ref{chap:4} hình thức hóa bài toán, Chương~\ref{chap:5} mô tả
chương trình và giao thức đánh giá, Chương~\ref{chap:6} trình bày kết quả, Chương~\ref{chap:7} trả lời mười câu hỏi,
Chương~\ref{chap:8} là hệ mua sắm đa tác tử, Chương~\ref{chap:9} thảo luận và kết luận. Mọi con số thực nghiệm trong báo
cáo được sinh tự động từ tệp JSON do chương trình ghi; mọi đoạn mã được trích nguyên văn từ mã đã chạy.

\begin{phathien}{Kết luận chính}
Ở cấu hình của đề bài ($\delta=20\%$), chính sách tối ưu khi đối thủ là tác tử luật là \emph{không bao giờ giảm giá}, đạt
$\SL{ch4_model/d20/optimal}$ mỗi bước, và tác tử luật vốn đã làm đúng như vậy. Q-learning học được chính sách tối ưu ở
$\SL{ch6_e/agg/n_opt}/10$ seed; chênh lệch so cặp với luật là $\SL{ch6_e/agg/diff_mean}$ với khoảng tin cậy $95\%$
$[\SL{ch6_e/agg/ci_lo};\,\SL{ch6_e/agg/ci_hi}]$, chứa $0$. Khi giảm giá rẻ hơn ($\delta=5\%$) luật không còn tối ưu và RL
vượt luật $\SL{ch6_f/005/diff_mean}$ mỗi bước. Tác tử dựa trên mô hình đề xuất ở câu~10 đạt tối ưu ở cả năm kịch bản thử.
\end{phathien}
